% ============================================================
%  Introduzione alla Parte III: Meccanica della trave rigida
% ============================================================

\chapter*{Introduzione Parte~III}
\addcontentsline{toc}{chapter}{Introduzione alla Parte~III}

La Parte~II ha sviluppato la meccanica dei sistemi di corpi rigidi nella sua forma più
generale, senza alcuna ipotesi sulla forma dei corpi. Questa parte ne è la
specializzazione geometrica: i corpi rigidi assumono la forma di \emph{travi}, corpi
caratterizzati da una dimensione --- la lunghezza --- predominante rispetto alle
dimensioni trasversali; la trattazione è riferita al caso delle \emph{travi piane},
con asse rettilineo contenuto in un piano e carichi e vincoli appartenenti allo stesso
piano, in continuità con la convenzione di lavoro seguita nel resto del volume.
La riduzione geometrica non è banale: la presenza di una dimensione privilegiata
introduce un riferimento intrinseco, individua componenti naturali degli spostamenti e
delle forze, e permette di descrivere lo stato interno del corpo mediante tre grandezze
scalari --- forza normale, taglio e momento flettente --- che non compaiono nella
meccanica del corpo rigido generico. Sono proprio queste grandezze a costituire il
ponte verso la meccanica della trave \emph{deformabile}, oggetto del Volume~II.

I risultati della Parte~II si trasferiscono senza modifiche: la struttura duale dei
problemi cinematico e statico, la classificazione dei sistemi vincolati, i metodi di
soluzione, la triade meccanica. Ciò che cambia è il contesto in cui questa struttura
opera. La specializzazione geometrica rende concreta la teoria generale, la arricchisce
di strumenti propri --- i diagrammi delle caratteristiche della sollecitazione, il
metodo per equivalenza, il catalogo delle tipologie strutturali, le linee di
influenza --- e produce quella varietà di sistemi che la pratica ingegneristica
incontra quotidianamente: travi ramificate, travi di Gerber, travi continue,
sistemi reticolari, telai.

\section*{Il programma della parte}

Il \textbf{Capitolo~\ref{cap:Trave_Rigida}} introduce la \emph{trave come corpo rigido}.
Si definisce il riferimento intrinseco, con l'asse $z$ lungo la linea d'asse e l'asse
$y$ nel piano della trave, e si fissa la convenzione grafica --- il tratteggio alla
fibra inferiore --- che orienta univocamente il riferimento locale nei sistemi di
travi. Il nucleo del capitolo è la specializzazione del vincolo di rigidità al
corpo-trave: le distorsioni --- longitudinale, trasversale e rotazionale nel riferimento
intrinseco --- modellano le discontinuità cinematiche imposte nelle sezioni interne, e
le reazioni duali del vincolo di rigidità si identificano con le caratteristiche della
sollecitazione: forza normale $N$, taglio $T$ e momento flettente $M$. Si introduce
il metodo per equivalenza per il calcolo della distribuzione completa di $N$, $T$, $M$
lungo la trave, e se ne ricava la forma differenziale --- le equazioni indefinite di
equilibrio. Si distingue infine la trave monoconnessa dalla trave pluriconnessa,
internamente iperstatica, e si descrivono gli svincoli interni che rilasciano
parzialmente il vincolo di continuità.

Il \textbf{Capitolo~\ref{cap:Diagrammi_CS}} affronta la \emph{costruzione dei diagrammi
	delle caratteristiche della sollecitazione} per la trave singola. Prima di procedere al
calcolo esplicito, si conduce un'analisi qualitativa preliminare fondata sulle equazioni
indefinite di equilibrio e sulle equazioni di salto: essa anticipa la forma dei
diagrammi --- uniforme, lineare, parabolica --- e la posizione dei punti di
discontinuità e di stazionarietà, fornendo uno strumento di verifica indipendente dal
calcolo. Le relazioni per aree --- che esprimono le variazioni di $N$, $T$, $M$ come
integrali dei carichi e del taglio --- completano il corredo di strumenti di lettura e
controllo. Per le travi iperstatiche si introduce il metodo del sistema principale, che
parametrizza la famiglia di soluzioni staticamente ammissibili nelle incognite
iperstatiche: la loro determinazione è rinviata al Volume~II, che dispone della legge
costitutiva della trave deformabile.

Il \textbf{Capitolo~\ref{cap:assemblaggi_travi}} estende l'analisi dalla trave singola
agli \emph{assemblaggi di travi}. Si mostra come la teoria dei sistemi di corpi rigidi
sviluppata nella Parte~II si trasferisca senza modifiche al contesto delle travi, e si
classificano le connessioni tra travi in base alle componenti di spostamento che
vincolano e alle caratteristiche della sollecitazione che trasmettono, ricavando per
ciascuna tipologia la molteplicità del collegamento. Si costruisce quindi un catalogo
delle principali tipologie strutturali, organizzato secondo la natura delle connessioni:
i \emph{sistemi a connessioni eterogenee} --- in cui coesistono connessioni di tipo
diverso --- e i \emph{sistemi a connessioni omogenee}, nelle due tipologie estreme dei
reticolari (connessioni esclusivamente a cerniera) e dei telai (connessioni
esclusivamente rigide). Per questa seconda famiglia si introduce la
\emph{doppia lettura topologica}: il sistema può essere visto con le travi come
protagoniste e le connessioni come vincoli, oppure con i nodi come protagonisti e
le travi come vincoli. Le due letture sono duali, conducono alla stessa formula di
classificazione e preparano i due approcci --- continuo e nodale --- che saranno
sviluppati nel Volume~II.

Il \textbf{Capitolo~\ref{cap:sistemi_conn_eterogenee}} ha carattere \emph{operativo}:
sviluppa la procedura per il tracciamento delle caratteristiche della sollecitazione
sui \emph{sistemi a connessioni eterogenee}, quelli in cui nodi rigidi, cerniere
interne e carrelli coesistono all'interno dello stesso sistema. Si introduce una
sequenza operativa in quattro passi --- equilibrio globale, isolamento dei
sottosistemi, scrittura delle caratteristiche della sollecitazione sui tronconi,
tracciamento sull'intera struttura --- e la si applica a quattro esempi paradigmatici:
la trave ramificata con nodi rigidi, l'arco (o portale) a tre cerniere, la trave di Gerber e la trave
continua. I primi tre esempi ammettono soluzione completa con il
solo equilibrio; le travi continue, iperstatiche, sono trattate fino alla
parametrizzazione della soluzione nelle incognite iperstatiche $\chi_j$, la cui
determinazione richiede le condizioni di congruenza elastica sviluppate nel Volume~II.

Il \textbf{Capitolo~\ref{cap:reticolari}} apre la trattazione dei sistemi a connessioni
omogenee, sviluppando in dettaglio l'analisi cinematica e statica dei \emph{sistemi
	reticolari piani}, a partire dall'inquadramento topologico e dalla formula di
Müller-Breslau. Sul versante cinematico si distingue la sovradeterminazione esterna --- dovuta a vincoli al suolo
sovrabbondanti --- da quella interna --- dovuta ad aste ridondanti --- e si studia
l'effetto delle distorsioni: quelle trasversali e rotazionali restano puramente locali
all'asta che le ospita, mentre le distorsioni longitudinali si propagano all'intera
travatura e il loro effetto si legge mediante il polo dell'asta, centro di rotazione
relativa dei due tronconi, che coincide con il polo di Ritter del problema statico
duale. Si introduce il principio di sostituzione asta--vincolo esterno, che consente
di decomporre travature complesse in sottostrutture più semplici preservando la
classificazione. Sul versante statico si sviluppano il metodo dei nodi --- con la
nozione di nodo canonico, la propagazione a cascata e i criteri di riconoscimento
immediato delle aste scariche --- e il metodo delle sezioni di Ritter, con l'analogia
della trave equivalente per le travature a correnti paralleli. La dualità tra polo
dell'asta e polo di Ritter percorre l'intero capitolo come filo conduttore, rendendo
visibile nella geometria discreta del reticolo la stessa simmetria cinematica-statica
che governa l'intera Parte~II.

Il \textbf{Capitolo~\ref{cap:telai_piani}} è il parallelo duale del precedente:
sviluppa l'analisi topologica e strutturale dei \emph{telai piani}, sistemi di travi
con connessioni esclusivamente rigide. Qui la doppia lettura topologica conduce alla
\emph{formula delle maglie} come misura dell'iperstaticità interna: ogni maglia chiusa
contribuisce con tre gradi di iperstaticità indipendentemente dai carichi, e questa
iperstaticità è \emph{intrinseca} alla tipologia --- non può essere eliminata, come
nei reticolari, rimuovendo aste ridondanti senza alterare la natura strutturale del
sistema. Si analizza la morfologia dei telai --- campate, piani, strutture miste con
svincoli interni --- e si sviluppano le equazioni esplicite della seconda lettura:
condizioni di compatibilità nodale ed equazioni di equilibrio nodale con $N$, $T$,
$M$ come incognite, nella struttura duale $\bm{B} = \bm{A}^\top$. Una sintesi
conclusiva unifica i tre casi fondamentali --- asta reticolare ($m=1$), trave con
svincolo ($m=2$), trave di telaio ($m=3$) --- in una progressione che nel Volume~II
diventerà la base del formalismo matriciale per sistemi di complessità arbitraria.
Il capitolo si chiude con la struttura parametrica della soluzione iperstatica,
illustrata sul portale doppiamente incastrato: le incognite iperstatiche $\chi_j$
parametrizzano la famiglia di soluzioni staticamente ammissibili, la loro
determinazione è rinviata al Volume~II.

Il \textbf{Capitolo~\ref{cap:TLV_travi}} chiude la Parte~III specializzando ai
sistemi di travi il principio dei lavori virtuali e la triade meccanica
congruenza--equilibrio--PLV introdotti in forma generale nel
Capitolo~\ref{cap:TLV}. Rispetto al contesto dei sistemi di corpi rigidi,
il vettore $\bm{s}$ dei dati cinematici si arricchisce: accoglie, oltre ai
cedimenti dei vincoli esterni e di collegamento, anche le \emph{distorsioni}
--- longitudinali, trasversali e rotazionali --- che le sezioni interne delle
travi possono subire. Parallelamente, il vettore $\bm{r}$ delle grandezze
reattive include, accanto alle reazioni vincolari esterne, le caratteristiche
della sollecitazione $N$, $T$, $M$ in una sezione interna, lette come reazioni
del vincolo di rigidità. Questa doppia estensione trasforma la triade in uno
strumento operativo di ampia portata, articolato in due formule duali per i
sistemi isostatici. Il \emph{metodo del carico unitario} fornisce il calcolo
diretto di una qualsiasi componente cinematica --- spostamento assoluto o
relativo, rotazione --- a partire dai dati cinematici assegnati, mediante la
soluzione di un problema statico fittizio con carico unitario duale. Il
\emph{principio di Müller-Breslau}, suo perfetto duale, fornisce il calcolo
di una qualsiasi componente reattiva --- esterna (reazione vincolare) o
interna (una delle caratteristiche della sollecitazione in una sezione) ---
a partire dai carichi applicati, mediante la soluzione di un problema
cinematico fittizio con cedimento o distorsione unitari nel vincolo duale.
Quando il carico reale è una forza di intensità unitaria che si sposta lungo
la struttura, la soluzione cinematica del problema duale descrive direttamente
la \emph{linea di influenza} della grandezza reattiva cercata: la curva che,
al variare della posizione del carico, ne fornisce l'ordinata. I due
strumenti --- carico unitario e Müller-Breslau --- sono la forma operativa
della dualità cinematica-statica che ha percorso l'intera Parte~III;
applicati ai sistemi iperstatici, costituiranno nel Volume~II la base del
\emph{metodo delle forze} e del \emph{metodo degli spostamenti}.

\section*{Generalità e specializzazione}

Un aspetto metodologico merita di essere dichiarato esplicitamente. I risultati
sviluppati in questa parte valgono per travi e sistemi di travi rigidi: la
deformabilità non è ancora introdotta, e le caratteristiche della sollecitazione
$N$, $T$, $M$ sono grandezze statiche interne il cui calcolo dipende solo
dall'equilibrio. Questa limitazione è al tempo stesso una forza: l'intera struttura
logica della Parte~III --- classificazione, metodi di soluzione, lettura dei
diagrammi --- è fondata sui soli principi della meccanica del corpo rigido, nella
loro forma specializzata alla geometria della trave.

La specializzazione ai sistemi a connessioni eterogenee, ai reticolari e 
ai telai, sviluppata negli ultimi tre capitoli, non aggiunge nuovi 
principi: applica la teoria generale a contesti geometrici che la 
pratica ingegneristica incontra con particolare frequenza. Il catalogo 
del Capitolo~\ref{cap:assemblaggi_travi} fissa il quadro generale --- 
la formula fondamentale $\nu = 3n_e - m$ e la mappa delle tipologie ---
e i capitoli successivi ne derivano le forme specializzate 
--- la formula di Müller-Breslau per i reticolari, la formula delle 
maglie per i telai --- sviluppando i metodi di soluzione e la 
morfologia propri di ciascuna tipologia. Nei sistemi a connessioni 
omogenee questa specializzazione produce una progressione unitaria --- 
dall'asta reticolare ($m=1$) alla trave con svincolo ($m=2$), fino alla 
trave di telaio ($m=3$) --- che nel Volume~II diventerà la struttura 
di un formalismo matriciale uniforme per sistemi di complessità 
arbitraria. Questa architettura --- teoria una volta per tutte, 
applicazioni per tipologia --- è la stessa che il Volume~II riproporrà 
per le strutture deformabili.


\section*{Il filo conduttore: specializzazione e dualità}

Il filo conduttore di questa parte è doppio. Il primo è la \emph{specializzazione
	progressiva}: dalla teoria generale dei corpi rigidi della Parte~II si scende alla
geometria della trave, poi alla trave singola, poi agli assemblaggi nelle loro diverse
tipologie. Ogni passo verso il particolare non aggiunge nuovi principi: applica quelli
già stabiliti a un contesto più concreto, arricchendoli degli strumenti propri della
specializzazione geometrica.

Il secondo è la \emph{dualità cinematica-statica}, che percorre la parte con la stessa
sistematicità della Parte~II. Le caratteristiche della sollecitazione sono duali alle
distorsioni del vincolo di rigidità. La doppia lettura topologica dei sistemi a
connessioni omogenee --- aste o nodi come protagonisti --- è la manifestazione più
visibile di questa dualità: in entrambi, reticolari e telai, essa produce equazioni
di struttura $\bm{B} = \bm{A}^\top$, rendendo visibile nella geometria discreta del
sistema la stessa simmetria cinematica-statica che governa l'intera Parte~II. Il
Capitolo~\ref{cap:TLV_travi} ne fornisce infine la forma operativa: il metodo del
carico unitario e il principio di Müller-Breslau sono due strumenti perfettamente
simmetrici --- l'uno calcola componenti cinematiche risolvendo un problema statico
fittizio, l'altro calcola grandezze reattive risolvendo un problema cinematico
fittizio --- che traducono la dualità in procedura di calcolo.

La motivazione più concreta per il Volume~II emerge, in modo duale, su entrambi i
versanti. Sul versante statico, l'iperstaticità interna dei sistemi più complessi
produce incognite che non possono essere determinate con la sola meccanica del corpo
rigido: le equazioni di equilibrio non bastano, e occorrono le equazioni costitutive
della trave deformabile. Sul versante cinematico, il duale è altrettanto stringente:
nei sistemi sovradeterminati, cedimenti vincolari e distorsioni genericamente assegnati
producono \emph{incompatibilità cinematiche} --- configurazioni che un sistema di corpi
rigidi non può soddisfare senza violare il vincolo di rigidità. Anche qui la
rigidità non basta: solo la deformabilità consente di accomodare quegli spostamenti
imposti che la meccanica del corpo rigido deve dichiarare impossibili. Iperstaticità
statica e impossibilità cinematica sono dunque due facce della stessa medaglia, e
insieme costituiscono la motivazione più solida per l'apertura del Volume~II.